\subsection{Polynominterpolation}
\label{subsec:real_functions_interpolation}
Bei der \emph{Polynominterpolation} \EnglishTerm{polynomial interpolation}
suchen wir ein Polynom, dessen Graph durch alle vorgegebenen Punkte verläuft.
Dafür verallgemeinern wir die Konstruktion einer Geraden durch zwei Punkte:
\par\medskip
{\centering
  \CachedTikzFigure{content/resources/chapters/real_functions/tikz/functions_interpolation_generalization.tex}\par
  \nopagebreak\smallskip
  % Keep the reference outside the cached figure so its number stays current.
  \makebox[17.2cm][r]{\textcolor{black!65}{%
    \textsuperscript{*}\,genauer in Satz~\ref{thm:real_functions_interpolation_unique}}}\par}
\medskip
\begin{TheoremBox}{Eindeutigkeit der Polynominterpolation}{real_functions_interpolation_unique}
  Ein Polynom vom Grad $n$, dessen Graph durch $n+1$ Punkte
  $P_i=(x_i,y_i)$ (paarweise verschiedenen $x_i$) verläuft, ist
  durch diese Punkte eindeutig bestimmt.
  \par\medskip
  {\centering
    \CachedTikzFigure{content/resources/chapters/real_functions/tikz/functions_interpolation_points.tex}\par}
\end{TheoremBox}
\begin{ExampleBox}{Parabel durch drei Punkte}{real_functions_interpolation_example}
  \begin{minipage}[t]{0.61\linewidth}
    \vspace{0pt}
    Gesucht ist ein Polynom vom Grad höchstens zwei durch
    $P=(0,0)$, $Q=(1,1)$ und $R=(2,0)$. Für die Koeffizienten verwenden wir
    den Ansatz
    \[
      p(x)=Ax^2+Bx+C,\qquad A,B,C\in\R.
    \]
    Das Einsetzen der Punkte in $p(x)$ führt zu drei linearen Gleichungen:
    \[
      \begin{aligned}
        Q(1,1):\quad p(1)=1
          &\Longleftrightarrow A\cdot1^2+B\cdot1+C=1\\[0.25em]
        R(2,0):\quad p(2)=0
          &\Longleftrightarrow A\cdot2^2+B\cdot2+C=0\\[0.25em]
        P(0,0):\quad p(0)=0
          &\Longleftrightarrow A\cdot0^2+B\cdot0+C=0
      \end{aligned}
    \]
  \end{minipage}\hfill
  \begin{minipage}[t]{0.36\linewidth}
    \vspace{0pt}\centering
    \CachedTikzFigure{content/resources/chapters/real_functions/tikz/functions_interpolation.tex}
  \end{minipage}

  \par\smallskip\noindent
  Diese bilden ein $3\times3$-LGS für die Unbekannte $A$, $B$ und $C$. Wir schreiben das LGS als
  erweiterte Koeffizientenmatrix:
  \[
    \begin{aligned}
      A+B+C&=1,\\
      4A+2B+C&=0,\\
      C&=0
    \end{aligned}
    \quad\Longleftrightarrow\quad
    \left[\begin{array}{@{}rrr|r@{}}
      1&1&1&1\\
      4&2&1&0\\
      0&0&1&0
    \end{array}\right].
  \]
  Dieses LGS lösen wir mit dem Gauß-Verfahren.
  \begin{ClozeArea}[34]{19}
    \textbf{Vorwärtselimination:} Die Zeilenoperation steht rechts neben der
    Matrix, auf die sie angewendet wird:
    \[
      \begin{array}{r@{\;}r@{\qquad}l}
        &\left[\begin{array}{@{}rrr|r@{}}
          1&1&1&1\\
          4&2&1&0\\
          0&0&1&0
        \end{array}\right]
        &\begin{array}{@{}l@{}}
          {}\\
          \smash{\TransformOp{Z_2-4Z_1}}\\
          {}
        \end{array}\\[1.0em]
        \Longleftrightarrow
        &\left[\begin{array}{@{}rrr|r@{}}
          1&1&1&1\\
          \textcolor{THMBlue}{0}&-2&-3&-4\\
          \textcolor{THMBlue}{0}&\textcolor{THMBlue}{0}&1&0
        \end{array}\right]
        &\text{Zeilenstufenform}
      \end{array}
    \]
    \textbf{Rückwärtseinsetzen:} Wir bestimmen die Unbekannten von unten nach oben:
    \begin{itemize}[leftmargin=1.4em,itemsep=0.55em,topsep=0.3em]
      \item Zeile 3: $\qquad$
        $C=0.$
      \item Zeile 2: $\qquad$
        $-2B-3C=-4
        \quad \overset{C=0}{\Longleftrightarrow} \quad
        -2B=-4
        \quad \Longleftrightarrow \quad
        B=2.$
      \item Zeile 1: $\qquad$
        $A+B+C=1
        \quad \overset{B=2,\,C=0}{\Longleftrightarrow} \quad
         A+2+0=1
         \quad \Longleftrightarrow \quad
         A=-1.$
    \end{itemize}
  \end{ClozeArea}
  Das gesuchte Polynom lautet
  \[
    p(x)=\Cloze{5.0cm}{-x^2+2x}.
  \]
\end{ExampleBox}
